\chapter*{The Appearances of the Risen Christ}


Now let's turn to the chronologically last, but by far the most important knot of events---the appearances Christ made following his resurrection. McDowell makes the following rational arguments in favor of their reality:

1. One could hardly suppose that any hostile, or at least critically inclined witnesses (of whom there was surely a decent number) would have missed any opportunity to expose the absurd rumors of a dead man appearing to the living. As of yet, however, no traces of any such denunciations have ever been found in any archive.

2. The witnesses to these appearances numbered in the dozens, and on one occasion (on the mountain in Galilee) he appeared to ``500 brethren'' at once. Is it conceivable to suppose that 500 people were simultaneously mistaken, or that all of them, to a man, conspired to lie about it?

3. People met the risen Christ both one-on-one and in groups, in a wide variety of emotional states, and at different times of the day (Mary Magdalene at dawn, the travelers on the road to Emmaus in the afternoon, and the Apostles in the evening, after dark). The diverse set of circumstances in which the appearances took place allows us to reject the commonly cited explanation of atheists that they were the product of ``hallucinations.''

4. McDowell draws attention to the fact that Christ's first appearances were not to his disciples, but to women---namely, Mary Magdalene and the Myrrhbearers. In his opinion, this is an important (albeit indirect) argument against the possibility of fraud. The issue here is that women's testimony held no legal weight according to Jewish law, so it would have made absolutely no sense to stage a hoax for them to witness.

As you can see, the persuasiveness of these arguments varies quite widely.

1. McDowell forgive me, but his point regarding the lack of recorded denunciations of the Resurrection is a paraphrase of a popular joke from the time of the ``Anti-Cosmopolitan Campaign.'' Western scientists find a piece of copper wire in an Egyptian tomb from the third millennium BCE; on that basis, they conclude that the ancient Egyptians were already using the telegraph. In response, Soviet scientists announce that no one has ever found any copper wires in Russian tombs from the third millennium BCE; therefore, the Russians of that time were already using the \emph{wireless} telegraph.

But on a serious note, the complete uniformity of archival data on a particular issue is generally a double-edged sword: in such cases, everything hinges on the original premise. For McDowell, who was raised in a democratic society, such uniformity comes off as a decisive argument in its favor; as for me, however (a product of Soviet totalitarianism), I actually see this as evidence of a targeted purge of the archives---an everyday occurrence! As an aside: if there's anything that convinces me of the historical authenticity of the Gospel texts and the lack of revisions made to them later on, it's precisely the discrepancies and inconsistencies they contain.\footnote{While I came to this conclusion completely independently, I was even more intrigued to find out later that Dmitry Merezhkovsky had already made a similar argument in his aforementioned apologia \emph{Jesus the Unknown}. (K.E.)}

2. If we rank Christ's known appearances by the strength of their supporting evidence, then the aforementioned ``appearance to the 500 brethren'' should really be at the very bottom of the list. Any student of psychology (or carnival barker, for that matter) can confirm that a crowd of that size can be convinced of just about anything in a pinch---unlike any of its members individually. In social psychology, this is known as ``the strengthening of a suggestive effect in group conditions,'' or in layman's terms, ``mob mentality''; a prime example of this is the ever-memorable seances of Anatoly Kashpirovsky.\footnote{A popular Russian psychic known for his large-scale demonstrations of his supposed powers. (Z.B.)}

These considerations, incidentally, apply just as strongly to another event that McDowell makes no mention of (and is technically not the subject of our analysis): the Ascension of the Lord. For what it's worth, note that this episode only appears in the Gospel of Luke (Lk 24:50-52) and Acts of the Apostles (Ac 1:2-11), which is based on the testimony of St. Paul, who did not witness it directly. Neither Matthew, Mark the Evangelist (who records Peter's account), nor, strangely enough, John---in other words, none of the direct participants in the events---say a single word about the Ascension. Returning to the ``appearance on the mountain in Galilee,'' I'd like to note that factually speaking, things are nowhere near as clear-cut as McDowell would have them (see below).

Arguments 3 and 4, on the contrary, seem quite reasonable to me. So now, keeping McDowell's aforementioned considerations in mind, let's take a closer look at all the times the risen Christ appeared to people who knew him closely. It is quite obvious that the testimony of this group of witnesses holds by far the most weight. We will not touch on Christ's appearances ``one-on-one'' to Peter the Apostle and James, Brother of Jesus, as such testimony could hardly be considered convincing in a legal sense, especially since the New Testament gives absolutely no details regarding these appearances. For the purposes of our analysis, our timeline of these appearances (just like McDowell, I consider this factor quite important) will follow Farrar's.

1. At dawn on the third day after the execution, when the guards discovered the empty tomb and headed off to inform the high priests, the Myrrhbearers arrived at the burial site. Upon entering the open crypt, they saw ``a young man dressed in white clothing'':

\begin{smallquote}
    And it happened, as they were greatly perplexed about this, that behold, two men stood by them in shining garments. Then, as they were afraid and bowed their faces to the earth, they said to them, ``Why do you seek the living among the dead? He is not here, but is risen! Remember how He spoke to you when He was still in Galilee, saying, `The Son of Man must be [...] crucified, and the third day rise again.''' \emph{And they remembered His words.}'' (Lk 24:4-8) \\
    \\
    But the angel answered, [...] ``And go quickly and tell His disciples that He is risen from the dead, and indeed \emph{He is going before you into Galilee; there you will see Him.}'' (Mt 28:5-7)
\end{smallquote}
Which the women did.

Thus, the Myrrhbearers' message regarding Christ's resurrection is what would legally be termed ``hearsay''; in this case, they are merely repeating the words of the ``men in shining garments.'' Furthermore, the men also have the women pass along a rather practical set of instructions to the disciples: hightail it out of Jerusalem on the double and lay low for a while in their homeland of Galilee. A wise idea, to say the least---over the next few days, things would be quite heated in Jerusalem, as the reaction of the angered and frightened high priests could be considerably harsh. In any case, the Sanhedrin would have a hard time getting their hands on the Apostles in Galilee, a territory that was traditionally hostile to Judeans.

There's one detail that's especially worth mentioning here. The Gospel of Matthew claims that Christ himself appeared to the women after their encounter with the angels. This statement (Mt 28:9-10), however, is quite strange for several reasons. The issue here is that the entire scene at the tomb (with the exception of this episode) is a rare case where the narratives of all four Evangelists, including John, line up almost completely; in the case of the Synoptic Gospels, these similarities extend even to the text itself. Could you really imagine that all the Evangelists (except Matthew) forgot to mention an event like this, whose significance needs no elaboration? What's more: the other Gospels effectively provide a direct refutation of Jesus' appearance to the Myrrhbearers (for example, Mk 16:9 and Lk 24:23).

It seems fairly obvious that the testimony of Peter (whose account forms the basis of the Gospel of Mark) and John is more credible in this case---for the simple reason that the two of them, unlike Matthew, were direct participants in the event (see below). It's also worth noting that in the latter's account, the only thing Christ says is ``Rejoice,'' following which he repeats, almost word for word, the angels' previous command to leave for Galilee (compare Mt 28:7 and 28:10). Therefore, we can assume that the frightened women simply took one of the ``shining angels'' for their teacher; Matthew, meanwhile, who was not a direct witness to the event, did no more than faithfully record their confused account.

2. In fact, the first to discover the empty tomb with the stone cast off was Mary Magdalene, who ran off at once to inform the Apostles. By the time she returned to the burial site, now with Peter and John in tow, the Myrrhbearers had already left. The Apostles examined the crypt, but found nothing more than empty burial shrouds, whereupon they shrugged their shoulders and went back to Jerusalem:

\begin{smallquote}
    But Mary stood outside by the tomb weeping [...] and she saw two angels in white sitting. [...] Then they said to her, ``Woman, why are you weeping?'' She said to them, ``Because they have taken away my Lord, and I do not know where they have laid Him.'' Now when she had said this, she turned around and saw Jesus standing there, \emph{and did not know that it was Jesus}. Jesus said to her, ``Woman, why are you weeping? Whom are you seeking?'' She, \emph{supposing Him to be the gardener}, said to Him, ``Sir, if You have carried Him away, tell me where You have laid Him, and I will take Him away.'' (Jn 20:11-15)
\end{smallquote}
And only after he makes several attempts at suggestion does she finally realize that the man speaking to her is Jesus himself.

A little strange not to immediately recognize one of your loved ones, don't you think? Okay, fine---grief, shock, the unsteady light of the early morning sun... But here's the interesting part. The angels in the white robes, who were talking with the Myrrhbearers just a short while earlier, disappeared the very moment Mary arrived at the tomb with Peter and John; the angels only turned up again once she was alone. And even when passing on the Teacher's instructions to his Apostles---that they should head to Galilee---the angels, for some reason, opted to use the women as a go-between instead of telling John and Peter directly, even though they were standing literally right there.

3. That same day, two of Christ's disciples, neither of whom were Apostles, were walking to the town of Emmaus:

\begin{smallquote}
    So it was, while they conversed and reasoned, that Jesus Himself drew near and went with them. \emph{But their eyes were restrained, so that they did not know Him.} (Lk 24:15-16)
\end{smallquote}
After carrying on a long (and theologically quite important) conversation, the three travelers reached Emmaus, where the disciples invited the stranger to share a meal with them. There, ``their eyes were opened, and \emph{in spite of his altered appearance}, they realized that the Lord was with them.''\footnote{Aleksandr Lopukhin, \emph{Biblical History in the Light of the Latest Research and Discoveries: the New Testament}, 1895, pp. 536-537.} Well, I'll be... What kind of ``altered appearance'' could he have had for his disciples not to recognize him as their teacher---even while conversing with him at length in broad daylight? What does Mark mean when he says that Christ ``appeared in another form'' (Mk 16:12)? How, strictly speaking, does it follow that the stranger was Christ? It's no wonder that the Apostles, who the two disciples proceeded to share their discovery with, were just as skeptical of their account as they were of the women's a short while earlier.

4. That same evening, ten of the Apostles were gathered inside a locked house to hide from the Judeans. Suddenly, Jesus appeared in the room with them and said, ```Peace be with you.' When He had said this, He showed them His hands and His side. Then the disciples were glad when they saw the Lord'' (Jn 20:19-20). The Lord was present with them bodily, but in an altered form; the Apostles even believed they were seeing a ghost. To convince them they were dealing with a living person, Jesus invited them to touch him, and concluded his demonstration by eating fish and honey with them. Conveniently absent from this meeting was Thomas---a skeptic who was constantly butting in with questions and doubts. Once he heard his comrades' story of Christ's appearance to them, he declared, ``Unless I [...] put my finger into the print of the nails, and put my hand into His side, I will not believe'' (Jn 20:25).

Looking back over this scene, it's hard to shake the strange impression that Jesus, in choosing to start his first meeting with the Apostles with a demonstration of his wounds, is using them as... let's say, a proof of identity. On the other hand, Thomas' reply also seems pretty strange when you think about it. Why would a normal, non-sadistic person suddenly get the urge to dig his finger into another man's wounds? For what it's worth, it was his reply in this particular episode that earned Thomas (``Twin'') his eventual nickname of ``Doubting Thomas,'' and not because he had any uniquely maniacal tendencies towards suspicion in general.

Clearly, something in his comrades' story set off a red flag, and I think I've got a pretty good guess of what it was. I'm sure he asked them, ``How was our teacher moving?'', and once he heard their answer, ``Normally, just like everyone else,'' he knew something was off. Since Christ's Heavenly Father resurrected him in the same body he was crucified in (as confirmed by the nature of his wounds), then that body, in theory, should have behaved accordingly. This raises a rather reasonable question: how genuine could those wounds really be if a man whose ankles had been pierced by nails the size of railroad spikes was walking around like nothing had ever happened?

Hardly any believer could bring himself to throw stones at Thomas---considering that the question of the nature of Christ's body following his resurrection immediately became one of the most pressing problems of theology. In any case, from the lengthy (and frankly, quite tangled) reasoning of St. Paul (1 Cor 15:35-54), we can conclude that it wasn't his directly resurrected earthly body; but if that's the case, why did it still have wounds? It's no surprise Thomas came to a simple, rather rustic solution to this question: ``I'll believe it when I feel it.''

There's another point here I just can't wrap my head around. Passing through walls and appearing inside a locked room is, without a doubt, a privilege reserved solely for ghosts. But there's no way a spirit can eat fish and honey---which means what we have before us is a being of the material world. It just doesn't add up. I personally have nothing against ghosts (especially those like the Ghost of Hamlet's Father or, say, the murdered samurai from \emph{Rashomon}\footnote{A Japanese film where multiple characters give contradictory accounts of a samurai's murder; the samurai himself posthumously testifies at the trial by way of a psychic. (Z.B.)}), but still, let's try to stay at least minimally consistent here and resist the temptation to change the rules partway into the game. I'm willing to admit that the Evangelist truthfully described his observations, but as for \emph{his interpretations of them}, I simply can't agree---at least until the two most likely hypotheses (according to ``Occam's Razor'') are refuted:

\begin{quote}
\begin{enumerate}
    \item[1.] The apparition was a wholly material being that lacked any actual supernatural traits.
    \item[2.] The apparition did not belong to the real world in any of its manifestations, and all of its traits (including its emphatically material ones) were equally illusory.
\end{enumerate}
\end{quote}

5. Thomas' wish was granted shortly afterwards, only eight days later. The Apostles, now fully assembled, were gathered once more inside a locked room. Once more, Christ appeared from inside the house, filling the Apostles with awe; after giving them another demonstration of the wounds in his hands and his side, he invited Thomas to touch them. The most interesting part, however, is what happened next. Christian commentators always write that Thomas actually \emph{felt} the Savior's wounds, following which he became a true believer---once and for all;\footnote{For example, see ``Thomas,'' \emph{Bible Encyclopedia}, vol. 4, p. 169. (K.E.)} here is Gladkov's retelling of the Gospel text:

\begin{smallquote}
    Suddenly the Lord stood in their midst, and after telling them, ``Peace be with you!'', He addressed Thomas: ``Take your finger here and see My hands.'' Thomas obeyed and felt the wounds of the nails in His hand with his finger. Then the Lord said to him: ``Take your finger and place it in My side, and do not be unbelieving.'' The Lord uncovered His pierced side, the wound in which was so great that one could fit one's hand inside it. Thomas, already convinced that the Lord's hands had been pierced with nails, now stretched his hand out towards the wound in His side, felt it, and falling before Him, exclaimed, ``My Lord and my God!''
\end{smallquote}

Now here's how it looked in reality:

\begin{smallquote}
    Jesus came, the doors being shut, and stood in the midst, and said, ``Peace to you!'' Then He said to Thomas, ``Reach your finger here, and look at My hands; and reach your hand here, and put it into My side. Do not be unbelieving, but believing.'' And Thomas \emph{answered and said to Him}, ``My Lord and my God!'' (Jn 20:26-28)
\end{smallquote}
That is to say, no ``feeling'' ever actually took place; Thomas acted the way any normal person would in his position, and his previous promises, as should have been expected, were merely a rhetorical device. All of this could be considered an insignificant detail, of course, if not for one ``but.'' Recall that Christ appeared before the Apostles ``in an altered form''; it was specifically his ability to appear from inside a locked house, as well as the nature of his wounds, that gave the Apostles any grounds to identify the man who appeared before them as Christ. But as for how genuine his wounds really were---as it turns out, no one knew.

6. Some time later, this time in Galilee, seven of the Apostles (which makes me wonder---where did the other four go?) went on a nighttime fishing trip on the Sea of Galilee:

\begin{smallquote}
    That night they caught nothing. But when the morning had now come, Jesus stood on the shore; \emph{yet the disciples did not know that it was Jesus.} (Jn 21:3-4)
\end{smallquote}
Following his suggestion, they cast their nets a second time---and this time, it was a smashing success:

\begin{smallquote}
    Therefore that disciple whom Jesus loved said to Peter, ``It is the Lord!'' Now when Simon Peter \emph{heard that it was the Lord}, he put on his outer garment [...] and plunged into the sea. But the other disciples came in the little boat. (Jn 21:7-8)
\end{smallquote}
Waiting for them on the shore was a campfire and a warm meal---baked fish and bread:

\begin{smallquote}
    Jesus said to them, ``Come and eat breakfast.'' \emph{Yet none of the disciples dared ask Him, ``Who are You?''---knowing that it was the Lord.} (Jn 21:12)
\end{smallquote}
This was followed by his famous dialogue with Peter: ``Tend my sheep.''

Does any of this make sense to you? Well, the fact they've failed to recognize their teacher yet again is only half the problem; at this point, I think we're more than used to it. In this case, however, we're talking about someone who---irrespective of their previous travels together---they have met twice in the past few days: they have conversed with him, ``felt'' his wounds, and shared a meal with him. I spent a long time trying to figure out what all of this reminded me of. Then suddenly, it hit me: this is how they teach orphans to call their adoptive father ``papa''...

7. Finally, last on Farrar's list is the ``appearance to the 500 disciples and Apostles on the mountain in Galilee.'' We have already discussed the inherent value of such collective testimony. Meanwhile, the factual side of the matter is as follows:

\begin{smallquote}
    Then the eleven disciples went away into Galilee, to the mountain which Jesus had appointed for them. When they saw Him, they worshiped Him; \emph{but some doubted.} (Mt 28:16-17)
\end{smallquote}
In short: business as usual.

So let's sum things up. If we appeal to the legal principles that McDowell holds so dear, we are forced to accept the following statement as fact: \emph{in none of the episodes we have examined would any court recognize that the risen Christ had been successfully identified by those who knew him closely.} This conclusion has, as I see it, two rather unexpected consequences:

1. If Christ's appearances really were the result of individual and collective hallucinations (an explanation highly favored by atheist commentators), then the witnesses would only have been presented with their own impressions of their teacher. In that case, each of them would have seen exactly (and only) what they wanted to see. Take St. Paul, for example (who at that time was still that savage persecutor of Christians, the Pharisee Saul): even though he had never seen Jesus in his life, he still immediately realized who he was talking to. However, the witnesses' clear (albeit somewhat hidden) doubts as to the authenticity of the risen Christ are incontrovertible proof that they were not presented with products of their own memories, but a real person of flesh and blood. Whether or not that man was Christ is another question altogether.

2. I've already taken the opportunity to note that, as I see it, the inconsistencies contained in the Gospel texts actually testify to the latter's authenticity. As it turns out, the Evangelists' mentions of the doubts expressed by the witnesses to these appearances are the clearest case of such a ``guarantee by contradiction.'' It is widely known that the Bodily Resurrection is the most crucial moment for Christianity in the entire story of Jesus of Nazareth:

\begin{smallquote}
    And if Christ is not risen, then our preaching is empty and your faith is also empty. (1 Cor 15:14)
\end{smallquote}
With this in mind, it stands to reason that it would be precisely this group of ``slip-ups'' that would vanish from the Gospels after the first targeted ``editing'' of the text. Any employee of Orwell's Ministry of Truth who missed a slip-up of this level would surely be vaporized on the spot.

Now let us recall McDowell's observation regarding the wide variety of circumstances in which these appearances took place (varying numbers of witnesses, varying times of day, etc.). In spite of this, let's try to find at least something linking all these events together. As it turns out, there really is such a common factor: \emph{poor lighting}. Appearances 2 and 6 took place in the early morning twilight, and appearances 4 and 5 took place during the evening---in a locked room, no less. Only appearance 3 (on the road to Emmaus) and perhaps appearance 7 (on the mountain in Galilee) took place in broad daylight, but note that it was during these appearances that the issue of identification was most problematic.

In classifying these events, it's easy to see that the two appearances to the Apostles in Jerusalem clearly stand alone. First, these are the only appearances where the apparition clearly and unambiguously identifies himself. Second, these are the only appearances where he carries clear traces of crucifixion---the characteristic set of wounds. It was obviously these two factors in combination that led witnesses to recognize the risen Christ with a somewhat greater degree of certainty in these appearances than in any of the others.

Moreover---take note!---it is these two appearances that hold the least meaning out of all of them: the Savior appears solely to show his wounds twice over, eat fish and honey, and admonish his disciples for their lack of faith. Conversely, there are only two appearances where Christ makes lengthy speeches on key points of religious dogma: on the road to Emmaus and on the Sea of Galilee. In these cases, as we may recall, the wounds are not present, nor is there a sufficiently high degree of confidence in the identity of the man making these speeches. Speaking of the wounds: they paint an interesting picture. So, he doesn't have them in Emmaus (in what universe could you assume the disciples didn't notice such a ``tiny detail'' in the hours they spent talking with him?), but then they appear the moment he makes his appearances in Jerusalem---yet by the time he makes his appearances in Galilee, they've vanished without a trace...

Now let's examine the timeline of the appearances from this perspective. The first witnesses are women---grief-stricken, frightened, but nonetheless singularly loyal to their deceased teacher; moreover, most of the explanation for what has happened comes from certain ``men in white robes.'' The next appearance is to two disciples who aren't members of Jesus' inner circle. And only after information regarding these encounters reaches the Apostles, and they become duly prepared for it, does Christ appear to them. But not all of them---it turns out the skeptic Thomas has had his housing rights revoked. And only after this stubborn man's comrades spend a week psychologically conditioning him does the next appearance follow, this time to all eleven Apostles. Moreover, it's hard to shake the impression that during this second appearance, Christ isn't interested in anyone but Thomas. You have to admit that this timeline follows a rather conspicuous pattern: at one end are the people who are the most excitable and suggestible, or know Christ the least, and at the other end are the people who most closely know him and are most capable of independent thought. In addition, each previous rung of the ladder has the ability to psychologically influence the next one.

And yet, when you go over the list of people Christ appeared to, it's missing a certain something---or rather, there's a completely inexplicable gaping hole. Jesus appeared to just about everybody: a pair of disciples from Emmaus who he casually knew, the ``500 brethren,'' and his clearly unfavored brother James. But there was one person he never once graced with a visit: his own mother.

This is so outrageous that Gladkov concludes his account of the appearances with this magnificent piece of reasoning:

\begin{smallquote}
    According to legend, Christ made His first appearance to Mary, the Mother of God. While the Evangelists say nothing about this appearance, \emph{it is hard to assume} [emphasis mine --- K.E.] that He would have appeared to the Apostles several times without once gratifying His mother with an appearance, considering how worried He was for Her in His death throes on the cross.
\end{smallquote}
It really is hard to assume such a thing, but ultimately, we are forced to. As for the words ``according to legend'': in this context, they sound---well, just like the magical introduction ``As we all know'' from the ever-memorable ``Announcements from TASS.''
